\section{本讲小结}
刚体平面运动由质心平动与绕质心转动组成，计算速度、动能和冲击都应保留这两部分。转动惯量必须说明转轴，移轴只适用于平行轴。平动、转动和绳长或滚动约束共同构成方程组。

摩擦方向由接触点相对速度决定，静摩擦大小由约束所需求出，再检验上限。进入纯滚动后要及时切换模型；离开轨道前应检查支持力与所需摩擦是否仍可实现。偏心冲击的恢复条件使用接触点速度，复摆周期则依赖悬轴惯量与质心力臂。
